\documentclass{article}
\usepackage{amsmath}
\title{Fibonacci Spiral: Mathematical Foundation}
\author{Mathematica}
\begin{document}
\maketitle

\section{Introduction}

The Fibonacci sequence \( F_n = F_{n-1} + F_{n-2} \) with \( F_0 = 0, F_1 = 1 \)
defines a spiral that converges to the golden ratio \( \phi = \frac{1 + \sqrt{5}}{2} \approx 1.618 \).

\section{Geometric Interpretation}

Each quarter-circle arc in the spiral has radius \( F_n \), producing the recurrence:
\[
  r_n = r_{n-1} + r_{n-2}
\]

The cumulative arc length approaches \( \phi \cdot r_n \) as \( n \to \infty \).

\section{Visualization Method}

The rendered animation plots 3D coordinates \( (x_n, y_n, z_n) \) derived from the
polar form \( r = e^{b\theta} \) where \( b = \frac{\ln \phi}{\pi/2} \).

\end{document}
